% ЗГЕНЕРОВАНО scripts/build_format_v2_chapter.py — не редагувати вручну.
% Джерело: books/ashtavakra-gita/poetychnyi/rozdil-9.md

\poemverse{9.1}{Зроблене й незроблене, \\ пари протилежностей — \\ коли вони втихли, \\ то чиї вони? \\ Знаючи це, будь тут, у відреченні, \\ схильним до відкидання — \\ і без обітниць.}
\poemverse{9.2}{У якогось щасливого, любий, \\ від споглядання того, \\ як поводиться світ, \\ бажання жити, \\ голод і жага пізнавати \\ дійшли до втишення.}
\poemverse{9.3}{Усе це нестале, \\ зіпсоване троїстою мукою, \\ без суті, \\ осуджене, гідне відкинення. \\ Так вирішивши, людина втихає.}
\poemverse{9.4}{Який же той час, вік або що, \\ коли в людей немає пар протилежностей? \\ Знехтувавши їх, \\ хто живе як випало — \\ той досяг би здійснення.}
\poemverse{9.5}{Побачивши, скільки різних поглядів \\ мають великі мудреці, \\ праведні йоги, — \\ хто, дійшовши відречення, \\ не втихне?}
\poemverse{9.6}{Хто, пізнавши образ свідомості, — \\ чи не вчитель? — \\ той через відречення, рівність і зʼєднаність \\ переправляє з сансари.}
\poemverse{9.7}{Дивись на видозміни сущого \\ як на самі лише елементи, \\ як є насправді, — \\ і тієї ж миті, вільний від пут, \\ станеш тим, хто стоїть у власній формі.}
\poemverse{9.8}{Сансара — це саме усталені схильності, так. \\ Відпусти їх усі. \\ Їх відкидання — через відкидання схильностей; \\ стояння в собі — нині, \\ так чи інакше.}
